\subsection{The unit ordinary terminal catalogue}\label{proof:Q5-C106}
\paragraph{Q106.1 Original model and what is proved}

Let (X, mu) be any nonempty finite probability space and let K be real with
\(K(x,\allowbreak y)=-K(y,\allowbreak x)\) and \(|K(x,\allowbreak y)|\le1\). All expectations and operators below use mu:

\fitdisplay{
Jv=(Ev)1,\quad A=J+K,\quad T=J-K=A^*,\quad
g=K1,\quad p=1+g,\quad q=1-g,\quad d=Eg^2\in[0,1/3],
}{}
\fitdisplay{
f_n=A^{n}\one,\qquad \bar f_n=T^{n}\one,\qquad h_n=Tf_{n-1}.
}{}

An \(f_n\) arm points entirely away from its root. A \(\bar f_n\) arm points entirely
toward its root. An \(h_n\) arm has its first edge toward the root and all
remaining edges toward its leaf. In particular \(f_1=p\) and \(h_1=q.\) 

At each root, pair equal-length original arms as far as possible using
\(f_\ell\) and \(h_\ell\) on each pair. Their actual product \(E_0\) satisfies
\(0<E_0\le 1\) after zero-mass atoms are removed. Entire pairs may instead be
reversed when the complete construction calls for it. The residual lengths
are distinct. The original equal pairs remain in the tree and edge ledger.

Write

\fitdisplay{
Q(d)=1-d/20000,\qquad
I_q(U)=E[q^{-1/3}|U|^{8/3}].
}{Q106.1.1}

Theorem A (complete ordinary terminal supply). Let Sigma be any
finite original terminal arm multiset with at least two distinct residual
arms after maximum equal pairing. For every positive original connector
length L, the tables in parts Q106.2--3 select a terminal orientation
\(F^{\rm ord}_{\Sigma,\allowbreak L}\), depending only on Sigma and L, such that

\fitdisplay{
\boxed{E[(E_rpT^L(E_tF^{ord}_{Sigma,L}))^8]\le Q(d).}
}{Q106.1.2}

Here \(E_t\) and \(E_r\) contain arbitrary genuine equal pairs at the terminal and
receiving root. The receiving root has an original outward unit arm,
whose factor is p. All L connector edges point toward the receiving root.
The statement covers residual size at least two; the same-direction PW4 and an f2-reference fourth for the chosen F are separate assertions.

Theorem B (root pair (1, a), all terminals and both lengths). At an
original internal root with residual exactly (1, a), \(a\ge 2,\allowbreak \) orient the unit
arm outward and the a-arm entirely inward. With any Sigma as in A and
any independent positive L, M, put

\fitdisplay{
Y=E_rpT^L(E_tF^{ord}_{Sigma,L}),\qquad
U=T^M(\bar f_aY).
}{}

Then

\fitdisplay{
\boxed{I_q(U)^3\le(1-d/72)Q(d)\le1.}
}{Q106.1.3}

All M outer edges point from the internal root toward the output root.
The \(\bar f_a\) factor is multiplied before \(T^M\); it is not commuted through T.

Theorem C (unit internal root, all terminals, outer length at least three).
If the internal root instead has residual exactly one outward unit arm,
the same Y satisfies, for every \(M\ge 3,\allowbreak \)

\fitdisplay{
\boxed{I_q(T^MY)^3\le(1-d)Q(d)\le1.}
}{Q106.1.4}

The applications of these ordinary roles are given in Sections~\ref{sec:five-chain} and \ref{sec:five-star-fork}.

\paragraph{Q106.2 Every even residual terminal, with every direction exception}

Write the distinct residual table as

\fitdisplay{
a_1<b_1<\cdots<a_j<b_j,\qquad j\ge1.
}{}

All rows below use every \(L\ge 1\) except where explicitly split by L. Each
table entry omits the common genuine \(E_t\) only. The displayed factor holds in the displayed range.


\begingroup\small
\setlength{\tabcolsep}{3pt}
\setlength{\LTleft}{0pt}\setlength{\LTright}{\fill}
\begin{longtable}{@{}p{0.2200\linewidth}p{0.2333\linewidth}p{0.2333\linewidth}p{0.2333\linewidth}@{}}
\caption{The unit ordinary terminal catalogue: parameter values.}\label{tab:parameters:q5-c106:1}\\
\toprule
Entry & Connector & F & Proof \\
\midrule\endfirsthead
\toprule
Entry & Connector & F & Proof \\
\midrule\endhead
\bottomrule\endfoot
\(j=1,\allowbreak  a_1=1,\allowbreak  b_1\ge 2\) & every L; & \(q f_{b_1}\); & Section~\ref{proof:Q5-C058-ordinary}. \\
\addlinespace[2pt]
\(j=1,\allowbreak  a_1=2,\allowbreak  b_1\ge 3\) & \(L=1\); & \(h_2 f_{b_1}\); & Sections~\ref{proof:Q5-C053-unit} and \ref{proof:Q5-C057}. \\
\addlinespace[2pt]
\(j=1,\allowbreak  (a_1,\allowbreak b_1)=(2,\allowbreak 3)\) & \(L=2\); & \(h_2 f_3\); & Section~\ref{proof:Q5-C085}. \\
\addlinespace[2pt]
\(j=1,\allowbreak  a_1=2,\allowbreak  b_1\ge 4\) & \(L=2\); & \(f_2 h_{b_1}\); & Section~\ref{proof:Q5-C085}. \\
\addlinespace[2pt]
\(j=1,\allowbreak  a_1=2,\allowbreak  b_1\ge 3\) & \(L\ge 3\); & \(h_2 f_{b_1}\); & Section~\ref{proof:Q5-C085}/Section~\ref{proof:Q5-C086} exact union. \\
\addlinespace[2pt]
\(j=1,\allowbreak  a_1=3,\allowbreak  b_1\ge 4\) & \(L=1\); & \(f_3 h_{b_1}\); & Section~\ref{proof:Q5-C089}. \\
\addlinespace[2pt]
\(j=1,\allowbreak  (a_1,\allowbreak b_1)=(3,\allowbreak 4)\) & \(L=2\); & \(h_3 f_4\); & Section~\ref{proof:Q5-C089}. \\
\addlinespace[2pt]
\(j=1,\allowbreak  a_1=3,\allowbreak  b_1\ge 5\) & \(L=2\); & \(f_3 h_{b_1}\); & Section~\ref{proof:Q5-C089}. \\
\addlinespace[2pt]
\(j=1,\allowbreak  a_1=3,\allowbreak  b_1\ge 4\) & \(L\ge 3\); & \(h_3 f_{b_1}\); & Section~\ref{proof:Q5-C089}. \\
\addlinespace[2pt]
\(j=1,\allowbreak  a_1\ge 4\) & every L; & \(h_{a_1}f_{b_1}\); & Section~\ref{proof:Q5-C090-ordinary}. \\
\addlinespace[2pt]
\(j\ge 2,\allowbreak  a_1=1\) & every L; & \(q f_{b_1} \prod_{i=2}^j h_{a_i}f_{b_i}\); & Section~\ref{proof:Q5-C092}. \\
\addlinespace[2pt]
\(j\ge 2,\allowbreak  a_1\ge 2\) & every L; & \(\prod_{i=1}^j h_{a_i}f_{b_i}\); & Section~\ref{proof:Q5-C090-ordinary}. \\
\end{longtable}
\endgroup


Section~\ref{proof:Q5-C058-ordinary} supplies 1-d/12000. The Section~\ref{proof:Q5-C085}/Section~\ref{proof:Q5-C086} union supplies Q(d). Section~\ref{proof:Q5-C089}
supplies 1-d/500. Section~\ref{proof:Q5-C090-ordinary} supplies exp(-d/6) for its multiple-pair row;
for \(a_1\ge 4,\allowbreak  j\ge 1\) its stronger expression is
exp(-(128j/5-24)d), whose coefficient is at least \(8/5.\) Section~\ref{proof:Q5-C092} supplies
\(\exp(-(1/4+16(j-2))d).\) Each is at most Q(d): for the exponential rows
use \(\exp(-ad)\le 1-ad/(1+ad),\allowbreak \) with \(d\le 1/3\) and the stated \(a\ge 1/6.\)
The latter fraction is at least d/7, which is more than d/20000.

Exhaustion is structural. A positive even residual has either \(j=1\) or
\(j\ge 2.\) With \(j=1\) the shorter length is 1, 2, 3 or at least 4. The cases
shorter 2 and 3 are partitioned by \(L=1,\allowbreak  L=2\) and \(L\ge 3,\allowbreak \) retaining the
actual shorter-pair exceptions. With \(j\ge 2\) the shorter length is 1 or
at least 2. These cases are disjoint and cover all even residual lists.
No direction is chosen using K, mu, or d.

\paragraph{Q106.3 Every odd residual terminal}

Write

\fitdisplay{
a<b_1<c_1<\cdots<b_k<c_k,\qquad k\ge1.
}{}

\paragraph{Q106.3.1 Minimum \(a\ge 2\)}

The catalogue in Section~\ref{proof:Q5-C097}, with the unit and adjacent proofs in Sections~\ref{proof:Q5-C070}, \ref{proof:Q5-C093} and \ref{proof:Q5-C095}, gives the following table:


\begingroup\small
\setlength{\tabcolsep}{3pt}
\setlength{\LTleft}{0pt}\setlength{\LTright}{\fill}
\begin{longtable}{@{}p{0.2200\linewidth}p{0.3500\linewidth}p{0.3500\linewidth}@{}}
\caption{The unit ordinary terminal catalogue: parameter values.}\label{tab:parameters:q5-c106:2}\\
\toprule
Entry & Connector & Actual F \\
\midrule\endfirsthead
\toprule
Entry & Connector & Actual F \\
\midrule\endhead
\bottomrule\endfoot
\(b_1\ge 5,\allowbreak  \text{any } a\ge 2,\allowbreak k\ge 1\) & every L; & \(\bar f_a \prod_{i=1}^k h_{b_i}f_{c_i}\). \\
\addlinespace[2pt]
\(b_1=4,\allowbreak  a=2\) or 3, any \(k\ge 1\) & every L; & \(\bar f_a \prod_{i=1}^k h_{b_i}f_{c_i}\). \\
\addlinespace[2pt]
\(a=2,\allowbreak b_1=3,\allowbreak k\ge 2\) & every L; & \(\bar f_2 f_3 h_{c_1} \prod_{i=2}^k h_{b_i}f_{c_i}\). \\
\addlinespace[2pt]
exact (2, 3, 4) & every L; & \(\bar f_2 f_3 h_4\). \\
\addlinespace[2pt]
exact (2, 3, c), \(c\ge 5\) & \(L=1\); & \(\bar f_2 h_3 f_c\). \\
\addlinespace[2pt]
exact (2, 3, c), \(c\ge 5\) & \(L\ge 2\); & \(\bar f_2 f_3 h_c\). \\
\end{longtable}
\endgroup


Each row gives at most \(\exp(-d/27)\le 1-d/28\le Q(d).\) Since \(b_1>a\ge 2,\allowbreak  b_1\) is either 3, 4 or at least 5; the value 3 forces \(a=2.\) The last
three rows split the remaining k and c cases. Thus all minima at least
two and all positive L are covered. In particular the \(L=1\) direction
for (2, 3, \(c\ge 5)\) is not changed to the \(L\ge 2\) direction.

\paragraph{Q106.3.2 Minimum \(a=1\)}

In the following table, later pairs mean the displayed increasing pairs,
and every row uses all \(L\ge 1\) unless a connector is specified.


\begingroup\small
\setlength{\tabcolsep}{3pt}
\setlength{\LTleft}{0pt}\setlength{\LTright}{\fill}
\begin{longtable}{@{}p{0.2200\linewidth}p{0.2333\linewidth}p{0.2333\linewidth}p{0.2333\linewidth}@{}}
\caption{The unit ordinary terminal catalogue: parameter values.}\label{tab:parameters:q5-c106:3}\\
\toprule
Entry & Connector & Actual F & Reference \\
\midrule\endfirsthead
\toprule
Entry & Connector & Actual F & Reference \\
\midrule\endhead
\bottomrule\endfoot
\(b_1\ge 4,\allowbreak  \text{any } k\ge 1\) & every L; & \(q \prod_{i=1}^k h_{b_i}f_{c_i}\); & Section~\ref{proof:Q5-C096-ordinary}. \\
\addlinespace[2pt]
\(b_1=3,\allowbreak k\ge 2\) & every L; & \(q f_3 h_{c_1} \prod_{i=2}^k h_{b_i}f_{c_i}\); & Section~\ref{proof:Q5-C096-ordinary}. \\
\addlinespace[2pt]
exact (1, 3, 4) & every L; & \(q f_3 h_4\); & Section~\ref{proof:Q5-C096-ordinary}. \\
\addlinespace[2pt]
exact (1, 3, c), \(c\ge 5\) & every L; & \(q f_3 h_c\); & Section~\ref{proof:Q5-C098}. \\
\addlinespace[2pt]
\(b_1=2,\allowbreak k\ge 3\) & every L; & \(q f_2 h_{c_1} \prod_{i=2}^k h_{b_i}f_{c_i}\); & Section~\ref{proof:Q5-C096-ordinary}. \\
\addlinespace[2pt]
exact (1, 2, c, b, e), \(4\le c<b<e\) & every L; & \(q f_2 h_c h_b f_e\); & Section~\ref{proof:Q5-C098}. \\
\addlinespace[2pt]
exact (1, 2, 3, b, e), \(5\le b<e\) & every L; & \(q f_2 h_3 h_b f_e\); & Section~\ref{proof:Q5-C098}. \\
\addlinespace[2pt]
exact (1, 2, 3, 4, e), \(e\ge 5\) & every L; & \(q f_2 h_e h_3 f_4\); & Section~\ref{proof:Q5-C098}. \\
\addlinespace[2pt]
exact (1, 2, c), \(c\ge 5\) & every L; & \(q f_2 h_c\); & Section~\ref{proof:Q5-C098}. \\
\addlinespace[2pt]
exact (1, 2, 4) & every L; & \(q f_2 h_4\); & Section~\ref{proof:Q5-C102}. \\
\addlinespace[2pt]
exact (1, 2, 3) & \(L=2\); & \(q f_2 \bar f_3\); & Section~\ref{proof:Q5-C102}. \\
\addlinespace[2pt]
exact (1, 2, 3) & \(L=1\) or \(L\ge 3\); & \(q f_2 h_3\); & Section~\ref{proof:Q5-C102}. \\
\end{longtable}
\endgroup


The Section~\ref{proof:Q5-C096-ordinary} rows give exp(-d), exp(-d/2), or 1-d/601, all at most Q(d).
Section~\ref{proof:Q5-C098} gives 1-d/601, 1-d/12000, exp(-d/2), or exp(-8d), as recorded
in its local proof. Section~\ref{proof:Q5-C102} gives 1-d/601 for the complete 123 table,
and 1-d/12000 for the complete 124 table. Thus every row supplies (Q106.1.2).

For exhaustion, \(b_1\) is 2, 3 or at least 4. When it is 3, k is 1 or
at least 2; \(k=1\) has \(c_1=4\) or \(c_1\ge 5.\) When \(b_1=2,\allowbreak \) first split \(k\ge 3,\allowbreak  k=2,\allowbreak \) and \(k=1.\) With \(k=2,\allowbreak \) either \(c_1\ge 4\); or \(c_1=3\) and \(b_2\ge 5\); or
\(c_1=3\), \(b_2=4\). With \(k=1,\allowbreak  c_1\) is 3, 4 or at least 5. All lengths in
these partitions are original integers. In the crossed five-arm row
each of the lengths 1, 2, 3, 4, e is used exactly once. There is no artificial
pair or omitted arm. The Section~\ref{proof:Q5-C102} \(L=2\) row uses \(\bar f_3,\allowbreak \) not \(h_3.\)

\paragraph{Q106.3.3 Genuine decorations, reversal, and degenerate kernels}

The moment statements include actual equal pairs. Alternatively
positivity gives \(E_rpT^L(E_tF)\le pT^LF,\allowbreak \) and hence (Q106.1.2) from the undecorated
row. This comparison retains all original pair edges in the orientation.
After removal of zero-mass atoms, p, \(q\ge \mu(x)>0,\allowbreak \) so inverse q powers below
are defined. If \(d=0,\allowbreak  g=0\) but K need not be zero; original leaf integration
gives every actual message one. All bounds then hold with equality.

Applying the same fixed rule to -K reverses every original arm, connector
and pair and exchanges A/T and p/q. This gives the full reversed version.


parts Q106.2 and 3 prove Theorem A from exact local statements.

\paragraph{Q106.4 The complete internal-root transfers}

The theorem of Section~\ref{proof:Q5-C066} applies to every real Y on the original space:

\fitdisplay{
I_q(T^M(\bar f_aY))^3\le(1-d/72)E|Y|^8,
\qquad a\ge2,\quad M\ge1.
}{Q106.4.1}

Substitute the actual Y from (Q106.1.2). This proves Theorem B with no
restriction on Sigma besides its at least two residual arms. The root's
original a-arm provides exactly \(\bar f_a,\allowbreak \) and its original unit arm provides
p. Thus (Q106.4.1) is applied in its actual operator order and with no added
reference arm. All original L and M are independent.

For Theorem C, the all-real map of Section~\ref{proof:Q5-C079} for every \(M\ge 3\) is

\fitdisplay{
I_q(T^MY)^3\le(1-d)E|Y|^8.
}{Q106.4.2}

Substituting (Q106.1.2) proves (Q106.1.4). The proof covers \(M\ge3\). For \(M=1,2\) the Section~\ref{proof:Q5-C079} costs need not leave enough loss in \(Q(d)\) for a contraction.

For a half-tree of type B put \(B_pair=(1-d/72)Q(d),\allowbreak \) and for type C put
\(B_unit=(1-d)Q(d).\) Both outputs satisfy \(I_q(U)^3\le B_type\le 1.\)

